\section{引言}

差分滤波器分别处理两个输入电压，再形成输出差值。两条通道的处理特性不一致时，共模输入会转化为差分输出。\citet{casas1996basics} 总结了两种共模抑制方法：匹配非耦合通道的传递函数，以及通过浮置公共节点耦合单端滤波器。本文研究采用对地支路的梯形网络，属于前一种配置。

需要匹配的电路量取决于测量频率和所接负载。\citet{hornero2021cm} 表明，元件失配可能使共模响应出现增益峰值和正相移。无源梯形网络的负载还会经输出阻抗取电流。因此，根据开路电压传递函数选择元件组合后，还需结合输出阻抗预测带载响应。

缓冲后的电阻传感接口中，样品电压与 $I(t)R(t)$ 成正比，参考电压与 $I(t)R_s$ 成正比。电阻在观测窗口内变化时，先滤波再取比会引入历史权重；通道响应失配时，电流变化还会影响常值电阻读数。因此，匹配规范需要同时确定线性网络的端接条件，以及动态比值实际估计的量。

本文计算指定端接下多级对地梯形网络可用于匹配的元件自由度，并将匹配条件与动态比值测量联系起来。利用经典梯形提取，将正值元件实现表示为系数坐标，进而分类讨论端接约束下的原始响应匹配与校准后匹配。精确构造和理想元件电路仿真核查负载依赖性；卷积分析区分常值电阻下的激励抑制、电阻历史权重，以及通道时延失配的一阶影响。

\begin{figure}[htbp]
\centering
\includegraphics[width=\linewidth]{figures/Fig_11_zh.pdf}
\caption{缓冲后比值模型的信号路径。高亮块为系数方程分析的无源网络，负载与探头导纳在其输出平面计入；传感电阻改变上游源幅值。}
\label{fig:system-plane}
\end{figure}

